\documentclass[a4paper,11pt]{article}
\usepackage[utf8]{inputenc}
\usepackage[T1]{fontenc}
\usepackage[french]{babel}
\usepackage[left=2cm,right=2cm,top=2cm,bottom=2cm]{geometry}
\usepackage{xcolor}
\usepackage{hyperref}
\usepackage{fontawesome5}
\usepackage{parskip}

% --- Couleurs Viveris ---
\definecolor{primary}{RGB}{0, 82, 147}
\hypersetup{colorlinks=true, urlcolor=primary}

\begin{document}

% --- EXPÉDITEUR ---
\noindent
\begin{minipage}[t]{0.5\textwidth}
    \textbf{Loïc Aron MBASSI EWOLO}\\
    Élève-Ingénieur Bac+5 -- Double diplôme ENSEA\\[3pt]
    \faMapMarker\ Cergy, France\\
    \faPhone\ (+33) 7 59 52 33 98\\
    \faEnvelope\ \href{mailto:loicaron.mbassiewolo@ensea.fr}{loicaron.mbassiewolo@ensea.fr}
\end{minipage}
\hfill
% --- DATE ---
\begin{minipage}[t]{0.4\textwidth}
    \raggedleft
    Cergy, le 5 février 2026
\end{minipage}

\vspace{1.2cm}

% --- DESTINATAIRE ---
\hfill
\begin{minipage}[t]{0.5\textwidth}
    \raggedleft
    \textbf{Viveris}\\
    Service Recrutement\\
    Massy (91), France
\end{minipage}

\vspace{1cm}

\noindent\textbf{Objet :} Candidature au stage Ingénieur Logiciel Embarqué -- Robot de Surveillance et d'Accueil

\vspace{0.8cm}

Madame, Monsieur,

Actuellement élève-ingénieur en double diplôme à l'ENSEA, spécialisé en systèmes embarqués et automatique, je suis à la recherche d'un stage de fin d'études à partir d'avril 2026. Le projet de robot autonome de surveillance développé par Viveris, combinant navigation sécurisée et interaction humaine, représente un défi technique passionnant que je souhaite relever à vos côtés.

L'approche « sûreté de fonctionnement » mentionnée dans votre offre fait directement écho à mon projet d'option à l'ENSEA : la conception d'un système de détection de défauts et de contrôle tolérant aux pannes pour drone. J'y ai modélisé le comportement d'un véhicule en conditions nominales et dégradées (perte d'actionneurs, biais capteurs), conçu des observateurs pour la détection de défauts (FDI), et implémenté des stratégies de reconfiguration de contrôleur (FTC). Cette expérience me permettra d'apporter une vision structurée à l'analyse des faiblesses des fonctions de déplacement actuelles et à l'élaboration de spécifications robustes.

Mon expérience en robotique autonome s'est consolidée dans le cadre du Challenge CoHoMa, où je développe des algorithmes SLAM et de fusion de capteurs (Lidar 3D, caméra de profondeur) sur Nvidia Jetson. La gestion des flux temps réel pour une navigation autonome sécurisée, les tests d'intégration sous Docker, et la programmation C/C++ embarquée sont autant de compétences directement transférables à votre projet.

Je maîtrise le cycle de développement complet que vous décrivez : de l'analyse des besoins aux tests de validation, en passant par les spécifications et les tests unitaires. Mon projet PROJET-TAXI (simulation multiprocessus en C) et mes expériences en laboratoire IoT m'ont appris la rigueur nécessaire au développement de systèmes critiques.

Je serais honoré de contribuer à l'amélioration des fonctions de déplacement de votre robot et de mettre mes compétences au service de Viveris. Je reste à votre disposition pour un entretien.

Je vous prie d'agréer, Madame, Monsieur, l'expression de mes salutations distinguées.

\vspace{0.8cm}

\textbf{Loïc Aron MBASSI EWOLO}\\
\textit{Élève-Ingénieur ENSEA / Polytechnique Yaoundé}

\end{document}
